% Insertion-ready note for the Vietnam application or its appendix.
% Requires the notation and bibliography definitions used in draft_of_13.tex.

\subsubsection{Numerical replication of the cubic-age CBPS specification}
\label{subsubsec:vietnam_cbps_replication}

The replication of Table~2 in \citet{sloczynski2025abadie} is exact at the
reported precision except for the CBPS-based Uysal estimator under the cubic
age specification. The published estimate and standard error are
\(0.208\) and \(0.232\), respectively, whereas the R implementation yields
\(0.210\) and \(0.233\). This difference does not originate from the sample,
the construction of the covariates, the definition of the normalized
estimator, or the variance formula. It arises entirely in the numerical
estimation of the CBPS propensity scores.

Both implementations target the sample analogue of the balancing condition
\begin{equation}
    \frac{1}{N}\sum_{i=1}^{N}
    \frac{Z_i-p(X_i;\alpha)}
         {p(X_i;\alpha)\{1-p(X_i;\alpha)\}}X_i=0,
    \label{eq:vietnam_cbps_numerical_moment}
\end{equation}
where \(X_i=(1,\text{age}_i,\text{age}_i^2,\text{age}_i^3)'\) in the cubic
specification. In the R workflow, \texttt{fit\_cbps\_alpha()} standardizes the
nonconstant columns of \(X\), initializes the coefficients at the logit
maximum-likelihood estimate, and applies Newton steps with backtracking until
the largest absolute sample moment in
\eqref{eq:vietnam_cbps_numerical_moment} is smaller than \(10^{-9}\). The
resulting propensity scores are passed to \texttt{tau\_u()}, which evaluates
the separately normalized Wald--IPW estimator in closed form. Finally,
\texttt{kappa\_analytic\_se\_one()} stacks the CBPS moments with the moments
defining the four normalized group means and \(\widehat\tau_u\), and computes
the robust sandwich standard error from this joint M-estimation system.

The authors' Stata replication instead calls \texttt{kappalate} with
\texttt{zmodel(cbps)}. In the currently installed \texttt{kappalate} version
1.1.0, the CBPS first step is implemented by Stata's \texttt{gmm} command.
For the cubic specification, the default Gauss--Newton algorithm does not
converge, while \texttt{technique(nr)} returns
\[
    \widehat\tau_u^{cb}=0.2083305,
    \qquad
    \widehat{\operatorname{se}}
    \bigl(\widehat\tau_u^{cb}\bigr)=0.2322452,
\]
which reproduces the published values after rounding. At this solution, the
largest absolute balancing moment is approximately \(5.45\times10^{-4}\).
The R solver instead reaches a maximum moment of approximately
\(4.41\times10^{-14}\) and produces
\[
    \widehat\tau_u^{cb}=0.2099954,
    \qquad
    \widehat{\operatorname{se}}
    \bigl(\widehat\tau_u^{cb}\bigr)=0.2327808.
\]
The R \texttt{CBPS} package with exact balancing gives a nearly
identical point estimate of \(0.2099755\), corroborating the tightly solved R
result.

The decomposition can be verified directly. Supplying the Stata propensity
scores to the R function \texttt{tau\_u()} reproduces \(0.2083305\), and
evaluating the R sandwich calculation at the Stata CBPS solution reproduces
the Stata standard error of \(0.2322452\). Hence the estimator and inference
functions agree conditional on the same propensity scores. The small table
discrepancy is solely a consequence of the numerical CBPS solution in a poorly
conditioned cubic polynomial specification, rather than a difference in the
econometric estimand or its implementation.

This distinction also matters for reproducibility. The published Stata file
installs the current \texttt{kappalate} release and does not pin either the
package version or the optimization technique. Its results can therefore
change when the package defaults change. A literal replication with version
1.1.0 requires \texttt{technique(nr)}. The main R analysis retains the more
tightly balanced solution; if the published value is reproduced separately,
the Stata version and optimization setting should be reported explicitly.
